\documentclass[11pt]{article}
\usepackage[margin=1in]{geometry}
\usepackage{amsmath,amssymb,amsthm,mathtools}
\usepackage{booktabs}
\usepackage{enumitem}
\usepackage[colorlinks=true,linkcolor=blue,citecolor=blue]{hyperref}

% ---------------------------------------------------------------
% Theorem environments
% ---------------------------------------------------------------
\newtheorem{theorem}{Theorem}
\newtheorem{lemma}{Lemma}
\newtheorem{proposition}{Proposition}
\newtheorem{corollary}{Corollary}
\theoremstyle{definition}
\newtheorem{definition}{Definition}
\newtheorem{assumption}{Assumption}
\newtheorem{designrule}{Rule}
\theoremstyle{remark}
\newtheorem*{remark}{Remark}

% ---------------------------------------------------------------
% Notation macros (one symbol, one meaning)
% ---------------------------------------------------------------
\newcommand{\R}{\mathbb{R}}
\newcommand{\C}{\mathbb{C}}
\newcommand{\Aset}{\mathcal{A}}          % data set (invariant set)
\newcommand{\Net}{\mathcal{N}_\theta}    % network
\newcommand{\Netstar}{\mathcal{N}^\star} % target of the network slot
\newcommand{\Meth}{\mathcal{S}}          % method
\newcommand{\Hard}{\mathcal{H}}          % hardness
\newcommand{\Err}{\mathcal{E}}           % error function
\newcommand{\Lip}{\operatorname{Lip}_{\Aset}}
\newcommand{\supA}[1]{\|#1\|_{\infty}}
\newcommand{\expt}[1]{\exp_{#1}}         % truncated exponential
\newcommand{\sys}{\mathrm{sys}}
\newcommand{\ran}{\mathrm{ran}}
\newcommand{\ex}{\mathrm{ex}}

\title{A hardness theory for learned time steppers\\[4pt]
\large Injected error of neural time-stepping methods of arbitrary order}
\author{Draft}
\date{}

\begin{document}
\maketitle

\begin{abstract}
This note gives a theory for the injected error of a learned time stepper.
A method is written as a known part plus a network slot.
The hardness of a method is the Lipschitz constant of the flow map minus the known part.
Under a capacity model, the injected error of a method is bounded by a power of its hardness.
For Taylor methods of order $p$ with exact lower slots, the hardness is of order $h^p$.
For the direct step ($p=0$), the hardness does not decrease with $h$.
Errors in lower slots add to the hardness if they are functions of the state, and add to an irreducible floor otherwise.
A test-mode calculus extends the results to unresolved modes, to methods that wrap the network in a Runge--Kutta integrator, and to methods that use past predicted states.
The rollout error is the product of an amplification factor, analyzed in \cite{ainslie2026}, and the injected error, analyzed here.
Proofs are in Appendix~\ref{app:proofs}.
\end{abstract}

\tableofcontents

% ===============================================================
\section{Notation and terms}\label{sec:notation}
% ===============================================================

Table~\ref{tab:symbols} lists the symbols.
Table~\ref{tab:terms} lists the defined terms.
Each term has one name and is used with that name throughout.

\begin{table}[htbp]
\centering
\caption{Symbols.}\label{tab:symbols}
\small
\begin{tabular}{@{}p{0.25\textwidth}p{0.70\textwidth}@{}}
\toprule
Symbol & Meaning \\
\midrule
$f$ & vector field of $\dot x = f(x)$, $x \in \R^{n_{\mathrm{x}}}$ \\
$g$ & $Df\, f$, the second time derivative as a function of the state \\
$\Phi_t$ & flow map of $f$ over time $t$ \\
$h$ & time step \\
$t_{\mathrm{end}}$ & rollout horizon; $J = t_{\mathrm{end}}/h$ is the number of rollout steps \\
$\Aset$ & data set, a compact invariant set of $\Phi_t$ \\
$d$ & intrinsic dimension of $\Aset$ \\
$x^{(k)}$ & $k$-th time derivative as a function of the state; $x^{(0)}(x) = x$ \\
$M_k$, $L_k$ & $\sup_{\Aset}\|x^{(k)}\|$ and $\Lip(x^{(k)})$; $L = L_1$ \\
$\Lip(u)$ & Lipschitz constant of $u$ on $\Aset$ \\
$\supA{u}$ & $\sup_{x\in\Aset}\|u(x)\|$ \\
$\Meth$ & a method \\
$G_\Meth$ & step map of the method $\Meth$ \\
$p$ & order of the network slot \\
$F_k$ & estimate in the lower slot of order $k$ \\
$e_k$, $e_k^\sys$, $e_k^\ran$ & slot error, its systematic part, its random part \\
$\varepsilon^{(k)}_\sys$, $\varepsilon^{(k)}_\ran$ & $\Lip(e^\sys_k)$ and $\sup\|e^\ran_k\|$ \\
$c_k$ & slot weight of the slot of order $k$ \\
$Q_h$, $Q_h^\ex$ & known part, and known part with exact lower slots \\
$\Net$ & network \\
$\Netstar$ & target of the network slot \\
$\Hard(\Meth)$ & hardness of the method $\Meth$ \\
$\delta(\Meth)$ & injected error of the method $\Meth$ \\
$\Err(B,\tau)$ & error function (Lipschitz bound $B$ of the fitted function, label noise $\tau$) \\
$W$, $m$ & number of network weights, number of effective training samples \\
$\nu$ & observation noise standard deviation \\
$\rho$, $\alpha$ & rate and exponent of the capacity model \\
$\delta_{\mathrm{opt}}$ & optimization floor \\
$\lambda$, $z$ & test-mode rate, $z = h\lambda$ \\
$\mu$, $\mu^\star$ & per-mode network value, per-mode value of the target \\
$R_\Meth(\mu;z)$ & method multiplier \\
$\psi(z)$ & known multiplier \\
$R$ & stability polynomial of a Runge--Kutta method \\
$\zeta$ & $h\mu^\star$ \\
$s$, $r$ & number of stages and order of a Runge--Kutta method \\
$a_{il}$, $b_i$, $\beta$, $\bar a$ & Runge--Kutta coefficients, $\beta = \sum_i|b_i|$, $\bar a = \max_i\sum_l|a_{il}|$ \\
$\kappa_i$ & Runge--Kutta stages \\
$\Lambda$ & Lipschitz bound of the trained network \\
$q$, $\phi_i$, $C_q$ & order, coefficients, and constant of a finite-difference estimate \\
$\nu'$ & bound on the noise amplitude \\
$\vartheta$ & scale factor in the homogeneity condition \\
$w$, $v$ & off-manifold variable $\dot x - f(x)$, velocity estimate \\
$\epsilon_j$ & rollout error after $j$ steps \\
$\chi$, $\Gamma$ & Lipschitz constant of the step map, amplification factor \\
$\xi$ & parasitic root \\
$\expt{n}(z)$ & truncated exponential $\sum_{k=0}^{n} z^k/k!$ \\
$\gamma$ & growth rate per unit time of the learned step map \\
\bottomrule
\end{tabular}
\end{table}

\begin{table}[htbp]
\centering
\caption{Defined terms.}\label{tab:terms}
\small
\begin{tabular}{@{}p{0.42\textwidth}p{0.22\textwidth}p{0.30\textwidth}@{}}
\toprule
Term & Where defined & Short meaning \\
\midrule
method, step map & Definition~\ref{def:method} & rule that maps $x_j$ to $x_{j+1}$ \\
lower slot, network slot & Definition~\ref{def:method} & places where estimates enter \\
exact slot, state-based slot, history-based slot & Definition~\ref{def:slottype} & types of lower slot \\
slot error, systematic part, random part & Definition~\ref{def:sloterror} & $e_k = e_k^\sys + e_k^\ran$ \\
affine method, known part, slot weight & Definition~\ref{def:affine} & $G_\Meth = Q_h + c_p\Net$ \\
Runge--Kutta wrapper, implicit Euler wrapper & Section~\ref{sec:nonaffine} & network inside integrator stages \\
exponential known part & Definition~\ref{def:testmode} & known multiplier $e^{z_{\mathrm{lin}}}$ \\
Taylor method of order $p$ & Definition~\ref{def:taylor} & known part is a Taylor sum \\
injected error & Definition~\ref{def:injected} & one-step error on $\Aset$ \\
error function, capacity model & Assumption~\ref{ass:capacity} & network error versus Lipschitz bound \\
hardness & Definitions~\ref{def:hardness}, \ref{def:hardnessNA} & Lipschitz constant of what the network supplies \\
target & Definition~\ref{def:hardnessNA} & $\Netstar$ \\
test mode, method multiplier, known multiplier & Definition~\ref{def:testmode} & linear scalar analysis \\
representable, per-mode hardness & Definition~\ref{def:permode} & existence and size of $\mu^\star$ \\
parasitic root & Definition~\ref{def:permode} & extra root from history-based slots \\
resolved mode, unresolved mode & Definition~\ref{def:testmode} & $|z|\le 1$, $|z|>1$ \\
rollout error, amplification factor & Section~\ref{sec:rollout} & error after $J$ steps \\
\bottomrule
\end{tabular}
\end{table}

% ===============================================================
\section{Setting}\label{sec:setting}
% ===============================================================

The system is $\dot x = f(x)$ with flow map $\Phi_t$.
The data are pairs $(x, \Phi_h(x))$ with $x$ drawn from the data set $\Aset$.

\begin{assumption}[Regularity]\label{ass:regular}
The data set $\Aset$ is compact and invariant, $\Phi_t(\Aset) = \Aset$ for all $t \in \R$.
For every order $k$ used below, $x^{(k)}$ is bounded and Lipschitz on $\Aset$.
The derivatives are defined by $x^{(0)}(x) = x$, $x^{(1)} = f$, and $x^{(k+1)} = D x^{(k)}\, f$.
\end{assumption}

For chaotic dissipative systems, $\Aset$ is the global attractor.
Its intrinsic dimension $d$ can be much smaller than $D$.

% ===============================================================
\section{Methods and slots}\label{sec:methods}
% ===============================================================

\begin{definition}[Method]\label{def:method}
A method $\Meth$ is a step map
\begin{equation}\label{eq:method}
x_{j+1} = G_\Meth\big(x_j;\, F_1,\dots,F_{p-1};\, \Net\big).
\end{equation}
The estimates $F_k \approx x^{(k)}$, $k = 1,\dots,p-1$, fill the lower slots.
The network $\Net$ fills the network slot.
The integer $p \ge 0$ is the order of the network slot.
\end{definition}

\begin{definition}[Slot types]\label{def:slottype}
A lower slot is an exact slot if $F_k = x^{(k)}$.
It is a state-based slot if $F_k$ is a function of the current state $x_j$ only.
It is a history-based slot if $F_k$ depends on predicted states $x_{j-1}, x_{j-2},\dots$.
\end{definition}

A finite-difference estimate is state-based during training, because it uses true states.
The same estimate is history-based during rollout, because it uses predicted states.

\begin{definition}[Slot error]\label{def:sloterror}
The slot error of the lower slot of order $k$ is $e_k = F_k - x^{(k)}$.
It is split as $e_k = e_k^\sys + e_k^\ran$.
The systematic part $e_k^\sys$ is a deterministic function of the state.
The random part $e_k^\ran$ has zero mean conditional on the state.
The sizes are $\varepsilon^{(k)}_\sys = \Lip(e_k^\sys)$ and $\varepsilon^{(k)}_\ran = \sup\|e_k^\ran\|$.
\end{definition}

\begin{definition}[Affine method]\label{def:affine}
A method is affine if
\begin{equation}\label{eq:affine}
G_\Meth = Q_h + c_p\,\Net, \qquad Q_h = \sum_{k=0}^{p-1} c_k F_k,
\end{equation}
with $F_0 = x^{(0)}$.
The map $Q_h$ is the known part.
The scalars $c_k$ are the slot weights.
The known part with exact lower slots is $Q_h^\ex = \sum_{k=0}^{p-1} c_k x^{(k)}$.
\end{definition}

\begin{definition}[Taylor method of order $p$]\label{def:taylor}
The Taylor method of order $p$ is the affine method with $c_k = h^k/k!$ for $k = 0,\dots,p$.
The Taylor method of order $0$ is the direct step, $G_\Meth = \Net$.
The Taylor method of order $1$ is the Euler step, $G_\Meth = x + h\Net$.
\end{definition}

\begin{definition}[Injected error]\label{def:injected}
The injected error of a method is $\delta(\Meth) = \supA{G_\Meth - \Phi_h}$.
\end{definition}

% ===============================================================
\section{Capacity model}\label{sec:capacity}
% ===============================================================

\begin{assumption}[Observation noise]\label{ass:noise}
The observed next state is $\Phi_h(x) + \eta$ with $\eta$ independent of $x$, zero mean, and standard deviation $\nu$.
The input state is observed without noise.
\end{assumption}

\begin{assumption}[Capacity model]\label{ass:capacity}
Consider a network class with $W$ weights and a final linear output layer.
Train it on $m$ effective samples to regress a function $u$ with $\Lip(u) \le B$ and label noise of standard deviation $\tau$.
The trained network $\hat u$ satisfies
\begin{equation}\label{eq:capacity}
\supA{\hat u - u} \le \Err(B, \tau) + \delta_{\mathrm{opt}}.
\end{equation}
The error function $\Err$ is nondecreasing in $B$ and jointly homogeneous of degree one,
\begin{equation}\label{eq:homog}
\Err(\vartheta B, \vartheta \tau) = \vartheta\, \Err(B, \tau) \quad \text{for all } \vartheta > 0.
\end{equation}
At fixed noise level $\nu$, the error function has the form
\begin{equation}\label{eq:rate}
\Err(B, \nu) = \rho(W, m, \nu)\, B^{\alpha}, \qquad 0 < \alpha \le 1.
\end{equation}
\end{assumption}

Two regimes give standard values of $\rho$ and $\alpha$ for Lipschitz functions on a set of dimension $d$.
In the approximation-limited regime, $\alpha = 1$ and $\rho \propto W^{-1/d}$ up to logarithmic factors \cite{yarotsky2017}.
In the data-limited regime, $\alpha = d/(d+2)$ and $\rho \propto \nu^{2/(d+2)} m^{-1/(d+2)}$ up to logarithmic factors \cite{stone1982}.
Both forms satisfy Eq.~\eqref{eq:homog}.
For samples taken from one trajectory, $m$ is the effective sample size after a mixing correction \cite{mohri2008}.
% VERIFY: the cited rates are stated for Euclidean domains and L2 error; the versions on a d-dimensional set and in sup norm need a citation.

\begin{lemma}[Output scaling does not change the injected error]\label{lem:scaling}
For an affine method with exact lower slots, under Assumptions~\ref{ass:noise} and~\ref{ass:capacity},
\begin{equation}\label{eq:scaling}
\delta(\Meth) \le \Err\big(\Lip(\Phi_h - Q_h^\ex),\, \nu\big) + \delta_{\mathrm{opt}}\,|c_p|.
\end{equation}
The bound depends on $c_p$ only through the optimization floor term.
\end{lemma}

Lemma~\ref{lem:scaling} is proved in Appendix~\ref{app:scaling}.
The injected error depends on what the known part leaves for the network.
It does not depend on how the network output is scaled.

% ===============================================================
\section{Hardness of affine methods}\label{sec:affine}
% ===============================================================

For an affine method with exact lower slots, the network is trained to fit $(\Phi_h - Q_h^\ex)/c_p$.

\begin{definition}[Hardness of an affine method]\label{def:hardness}
The hardness of an affine method is
\begin{equation}\label{eq:hardness}
\Hard(\Meth) = \Lip\big(\Phi_h - Q_h^\ex\big).
\end{equation}
\end{definition}

\begin{theorem}[Hardness bounds the injected error]\label{thm:affine}
Let $\Meth$ be an affine method whose lower slots are exact or state-based.
Under Assumptions~\ref{ass:noise} and~\ref{ass:capacity},
\begin{equation}\label{eq:thm1}
\delta(\Meth) \le \rho\,\Big[\Hard(\Meth) + \sum_{k=1}^{p-1} |c_k|\, \varepsilon^{(k)}_\sys\Big]^{\alpha}
+ \sum_{k=1}^{p-1} |c_k|\, \varepsilon^{(k)}_\ran + |c_p|\,\delta_{\mathrm{opt}},
\end{equation}
where $\rho = \rho(W, m, \nu_{\mathrm{eff}})$ and $\nu_{\mathrm{eff}}^2 = \nu^2 + \sum_{k} c_k^2\, \mathrm{Var}(e_k^\ran)$.
\end{theorem}

Theorem~\ref{thm:affine} is proved in Appendix~\ref{app:affine}.
The systematic part of a slot error adds to the hardness.
The random part of a slot error adds to a floor that no network can remove.

\begin{lemma}[Taylor known part]\label{lem:taylor}
Let $Q^{(n)}_h = \sum_{k=0}^{n} \frac{h^k}{k!}\, x^{(k)}$ with $n \ge 0$.
Under Assumption~\ref{ass:regular},
\begin{align}
\supA{\Phi_h - Q^{(n)}_h} &\le \frac{M_{n+1}\, h^{n+1}}{(n+1)!}, \label{eq:taylorsup}\\
\Lip\big(\Phi_h - Q^{(n)}_h\big) &\le \frac{L_{n+1}\, e^{hL}\, h^{n+1}}{(n+1)!}. \label{eq:taylorlip}
\end{align}
\end{lemma}

\begin{lemma}[Direct step]\label{lem:direct}
Under Assumption~\ref{ass:regular}, $\|\Phi_h(x) - \Phi_h(x')\| \ge e^{-hL}\|x - x'\|$ for all $x, x' \in \Aset$.
The hardness of the direct step satisfies $\Hard \ge e^{-hL}$.
\end{lemma}

Lemmas~\ref{lem:taylor} and~\ref{lem:direct} are proved in Appendices~\ref{app:taylor} and~\ref{app:direct}.

\begin{corollary}[Hardness of Taylor methods]\label{cor:taylor}
For the Taylor method of order $p$ with exact lower slots,
\begin{equation}\label{eq:cortaylor}
\Hard = \Lip\big(\Phi_h - Q^{(p-1)}_h\big) \le \frac{L_p\, e^{hL}\, h^p}{p!}, \qquad p \ge 1,
\end{equation}
and $\Hard \ge e^{-hL}$ for $p = 0$.
\end{corollary}

Corollary~\ref{cor:taylor} follows from Lemmas~\ref{lem:taylor} and~\ref{lem:direct} with $n = p-1$.
The Euler step has hardness at most $L e^{hL} h$.
The direct step has hardness at least $e^{-hL}$.
This is the identity argument.
The Taylor method of order $p$ removes one more power of $h$ for each exact lower slot.

\begin{designrule}[Order selection]\label{rule:order}
The bound in Eq.~\eqref{eq:cortaylor} for order $p$ is smaller than the bound for order $p-1$ if and only if
\begin{equation}\label{eq:rule1}
h\, L_p < p\, L_{p-1}.
\end{equation}
\end{designrule}

Rule~\ref{rule:order} follows from the ratio of the two bounds in Eq.~\eqref{eq:cortaylor}, which is $h L_p / (p L_{p-1})$.
For $p = 1$ it reads $h L < 1$, with $L_0 = 1$.

\begin{lemma}[Finite-difference slot error]\label{lem:fd}
Let $F_1(x) = h^{-1}\sum_{i=0}^{q} \phi_i\, \Phi_{-ih}(x)$ with coefficients that satisfy
\begin{equation}\label{eq:fdorder}
\sum_{i=0}^{q} \phi_i \frac{(-i)^k}{k!} = \begin{cases} 1 & k = 1,\\ 0 & k \ne 1,\end{cases} \qquad k = 0,\dots,q.
\end{equation}
Under Assumption~\ref{ass:regular}, the systematic part of the slot error satisfies
\begin{equation}\label{eq:fd}
\varepsilon^{(1)}_\sys = \Lip(F_1 - f) \le C_q\, L_{q+1}\, e^{qhL}\, h^{q}, \qquad C_q = \sum_{i=0}^{q} |\phi_i|\,\frac{i^{q+1}}{(q+1)!}.
\end{equation}
\end{lemma}

Lemma~\ref{lem:fd} is proved in Appendix~\ref{app:fd}.
The two-point backward difference has $q = 1$, $\phi_0 = 1$, $\phi_1 = -1$, and $C_1 = 1/2$.
The three-point backward difference has $q = 2$, $(\phi_0, \phi_1, \phi_2) = (3/2, -2, 1/2)$, and $C_2 = 1$.

\begin{designrule}[Acceptable lower slots]\label{rule:fd}
In a Taylor method of order $p$, a state-based lower slot of order $k$ keeps the hardness of order $h^p$ if and only if $\varepsilon^{(k)}_\sys = O(h^{p-k})$.
For a finite-difference estimate of $x^{(1)}$ of order $q$, this condition is $q \ge p - 1$.
\end{designrule}

Rule~\ref{rule:fd} follows from Eq.~\eqref{eq:thm1} with $c_k = h^k/k!$, and from Eq.~\eqref{eq:fd}, which gives $c_1 \varepsilon^{(1)}_\sys = O(h^{q+1})$.

\begin{corollary}[Noise floor of finite-difference slots]\label{cor:noise}
Under Assumption~\ref{ass:noise}, a finite-difference estimate from Lemma~\ref{lem:fd} evaluated on noisy states has random part with $\varepsilon^{(1)}_\ran \le \nu' \sum_i |\phi_i| / h$, where $\nu'$ bounds the noise amplitude.
The floor term $c_1 \varepsilon^{(1)}_\ran \le \nu' \sum_i |\phi_i|$ does not decrease with $h$.
The order-$p$ term of the flow map is larger than this floor only if
\begin{equation}\label{eq:noisewindow}
h > \Big(\frac{p!\, \nu' \sum_i |\phi_i|}{M_p}\Big)^{1/p}.
\end{equation}
\end{corollary}

Corollary~\ref{cor:noise} gives a lower limit on $h$.
Rule~\ref{rule:order} gives an upper limit on $h$.

% ===============================================================
\section{Non-affine methods}\label{sec:nonaffine}
% ===============================================================

A Runge--Kutta wrapper places the network inside stages.
The step map is
\begin{equation}\label{eq:rk}
G_\Meth[\Net](x) = x + h \sum_{i=1}^{s} b_i\, \kappa_i, \qquad \kappa_i = \Net\Big(x + h\sum_{l<i} a_{il}\, \kappa_l\Big),
\end{equation}
for an explicit method with $s$ stages.
A method of the form Eq.~\eqref{eq:rk} is a Runge--Kutta wrapper.
The implicit Euler wrapper is the method $x_{j+1} = x_j + h\Net(x_{j+1})$.
Set $\beta = \sum_i |b_i|$ and $\bar a = \max_i \sum_l |a_{il}|$.

\begin{assumption}[Target of the network slot]\label{ass:target}
There is a map $\Netstar$ with $G_\Meth[\Netstar] = \Phi_h$ on $\Aset$.
For a Runge--Kutta method of order $r$, $\Netstar = f + O(h^r)$ in $\Lip$.
\end{assumption}

For analytic $f$ and small $h$, the map $\Netstar$ in Assumption~\ref{ass:target} is the inverse modified vector field of \cite{zhu2020imde}.
The expansion fails when $h|\lambda| \gtrsim 1$ for an eigenvalue $\lambda$ of $Df$.
Section~\ref{sec:testmode} treats that case.

\begin{assumption}[Wrapper regression]\label{ass:wrapper}
Training $G_\Meth[\Net]$ on flow-map data gives a network whose distance to $\Netstar$ obeys Assumption~\ref{ass:capacity} with $u = \Netstar$ and label noise $\nu/h$.
\end{assumption}

\begin{lemma}[Runge--Kutta sensitivity]\label{lem:rk}
Let $\mathcal N_1$, $\mathcal N_2$ be defined on a set that contains all stage points, and let $\mathcal N_1$ be $\Lambda$-Lipschitz there.
For the explicit method in Eq.~\eqref{eq:rk},
\begin{equation}\label{eq:rksens}
\supA{G_\Meth[\mathcal N_1] - G_\Meth[\mathcal N_2]} \le h\,\beta\,(1 + h\Lambda\bar a)^{s-1}\, \|\mathcal N_1 - \mathcal N_2\|_{\infty}.
\end{equation}
\end{lemma}

\begin{definition}[Hardness of a non-affine method]\label{def:hardnessNA}
The map $\Netstar$ in Assumption~\ref{ass:target} is the target.
For an affine method with exact lower slots, the target is $(\Phi_h - Q_h^\ex)/c_p$.
The hardness of the method in Eq.~\eqref{eq:rk} is $\Hard(\Meth) = h\,\Lip(\Netstar)$.
\end{definition}

\begin{theorem}[Hardness of Runge--Kutta wrappers]\label{thm:nonaffine}
Under Assumptions~\ref{ass:capacity}, \ref{ass:target}, and~\ref{ass:wrapper}, with $\Lambda$ a Lipschitz bound for the trained network,
\begin{equation}\label{eq:thm2}
\delta(\Meth) \le \beta\,(1 + h\Lambda\bar a)^{s-1}\, \Big[\Err\big(\Hard(\Meth), \nu\big) + h\,\delta_{\mathrm{opt}}\Big].
\end{equation}
For a Runge--Kutta method of order $r$, $\Hard(\Meth) = h\big(L + O(h^r)\big)$.
\end{theorem}

Lemma~\ref{lem:rk} and Theorem~\ref{thm:nonaffine} are proved in Appendices~\ref{app:rk} and~\ref{app:nonaffine}.
For methods with $b_i \ge 0$, $\beta = 1$.
At small $h$, every Runge--Kutta wrapper has the same leading hardness $hL$ as the Euler step.
The order $r$ changes the target $\Netstar$ at order $h^r$.
It also changes the behavior of unresolved modes, as shown in Section~\ref{sec:testmode}.

% ===============================================================
\section{Test-mode calculus}\label{sec:testmode}
% ===============================================================

Theorems~\ref{thm:affine} and~\ref{thm:nonaffine} use expansions in $hL$.
This section applies each method to one linear mode.
The analysis holds for any value of $h\lambda$.

\begin{definition}[Test mode]\label{def:testmode}
The test mode is $\dot x = \lambda x$ with $\lambda \in \C$ and $z = h\lambda$.
In the test mode, the network slot is $\Net(x) = \mu x$ with $\mu \in \C$.
The method multiplier $R_\Meth(\mu; z)$ is defined by $G_\Meth(x) = R_\Meth(\mu; z)\, x$ when all history-based slots use the true trajectory.
For an affine method, $R_\Meth(\mu;z) = \psi(z) + c_p \mu$, where $\psi(z)$ is the known multiplier.
A mode is a resolved mode if $|z| \le 1$ and an unresolved mode if $|z| > 1$.
An exponential known part is a known part with known multiplier $e^{z_{\mathrm{lin}}}$, where $z_{\mathrm{lin}} = h\lambda_{\mathrm{lin}}$ and $\lambda_{\mathrm{lin}}$ is the rate of the exact linear operator on the mode.
\end{definition}

\begin{definition}[Representable, per-mode hardness, parasitic root]\label{def:permode}
The test mode is representable by $\Meth$ if some $\mu^\star$ satisfies $R_\Meth(\mu^\star; z) = e^z$.
If it is not representable, the method has the irreducible per-step error $\min_\mu |R_\Meth(\mu; z) - e^z|$.
The per-mode hardness is
\begin{equation}\label{eq:permode}
\Hard(\Meth; z) = \big|\partial_\mu R_\Meth(\mu^\star; z)\; \mu^\star\big|.
\end{equation}
A parasitic root is a root of the characteristic polynomial of the rollout recurrence other than $e^z$.
It exists only for methods with history-based slots.
\end{definition}

For an affine method, $\Hard(\Meth; z) = |e^z - \psi(z)|$.
This equals Definition~\ref{def:hardness} applied to the test mode, because $\Phi_h - Q^\ex_h$ is then the map $x \mapsto (e^z - \psi(z))x$.
For the method in Eq.~\eqref{eq:rk}, $\Hard(\Meth; z) = |\zeta R'(\zeta)|$ with $\zeta = h\mu^\star$.
The factor $|\zeta|$ equals Definition~\ref{def:hardnessNA} applied to the test mode.
The factor $|R'(\zeta)|$ is the exact test-mode value of the sensitivity bound in Eq.~\eqref{eq:rksens}.

\begin{proposition}[Per-mode hardness of standard methods]\label{prop:table}
Table~\ref{tab:permode} gives the per-mode hardness, the representability condition, and the parasitic roots.
\end{proposition}

\begin{table}[htbp]
\centering
\caption{Per-mode hardness. Here $\expt{n}(z) = \sum_{k=0}^n z^k/k!$ and $R$ is the stability polynomial of the Runge--Kutta method. The exponential known part uses the exact linear part $z_{\mathrm{lin}}$ of $z$.}\label{tab:permode}
\small
\begin{tabular}{@{}p{0.34\textwidth}p{0.27\textwidth}p{0.17\textwidth}p{0.12\textwidth}@{}}
\toprule
Method & $\Hard(\Meth; z)$ & Representable & Parasitic root \\
\midrule
Direct step & $|e^{z}|$ & all $z$ & none \\
Euler step & $|e^z - 1|$ & all $z$ & none \\
Implicit Euler wrapper & $|e^z|\,|e^z - 1|$ & all $z$ & none \\
Explicit Runge--Kutta wrapper & $|\zeta R'(\zeta)|$, $R(\zeta) = e^z$ & $e^z \in R(\C)$ & none \\
Taylor order $p$, exact lower slots & $|e^z - \expt{p-1}(z)|$ & all $z$ & none \\
Taylor order 2, backward difference, history-based & $|e^z + e^{-z} - 2|$ & all $z$ & $e^{-z}$ \\
Exponential known part & $|e^{z_{\mathrm{lin}}}|\,|e^{z - z_{\mathrm{lin}}} - 1|$ & all $z$ & none \\
\bottomrule
\end{tabular}
\end{table}

Proposition~\ref{prop:table} is proved in Appendix~\ref{app:table}.
For affine methods, the per-mode hardness is the distance between $e^z$ and the known multiplier.
The direct step has known multiplier $0$.
The Euler step has known multiplier $1$.
The Taylor method of order $p$ has known multiplier $\expt{p-1}(z)$.

\begin{proposition}[Crossovers on the real axis]\label{prop:cross}
For real $z$:
\begin{enumerate}[label=(\alph*)]
\item The Euler step has smaller per-mode hardness than the direct step if and only if $z > -\ln 2$.
\item The implicit Euler wrapper has smaller per-mode hardness than the Euler step if and only if $z < 0$.
\item The implicit Euler wrapper has smaller per-mode hardness than the direct step if and only if $z < \ln 2$.
\item The Taylor method of order 2 with exact lower slots has smaller per-mode hardness than the direct step if and only if $z > -1$.
\item The Taylor method of order 2 with exact lower slots has smaller per-mode hardness than the Euler step if and only if $z > z_\ast$, where $z_\ast \approx -1.5936$ solves $2e^{z} - 2 - z = 0$.
\item The Taylor method of order 2 with the history-based backward difference has smaller per-mode hardness than both the direct step and the Euler step if and only if $z > -\ln 2$.
\item As $z \to -\infty$, the Taylor method of order 2 with exact lower slots has per-mode hardness $|z| - 1 + e^z$, which grows without bound.
\end{enumerate}
\end{proposition}

\begin{proposition}[Imaginary axis]\label{prop:imag}
For $z = \mathrm{i}\omega$ with real $\omega$, the per-mode hardness of the direct step is $1$, of the Euler step is $2|\sin(\omega/2)|$, of the Taylor method of order 2 with exact lower slots is $|e^{\mathrm{i}\omega} - 1 - \mathrm{i}\omega| = \omega^2/2 + O(\omega^3)$, and of the Taylor method of order 2 with the history-based backward difference is $4\sin^2(\omega/2)$.
\end{proposition}

\begin{proposition}[Representability limit of explicit Runge--Kutta wrappers]\label{prop:floor}
For real $\zeta$, the stability polynomial of the classical second-order Runge--Kutta method has minimum $1/2$ at $\zeta = -1$.
The stability polynomial of the classical fourth-order Runge--Kutta method has minimum $0.2704$ at $\zeta \approx -1.5961$.
A decaying real test mode with $e^z$ below this minimum is not representable by a real $\mu$.
The limits are $z < -\ln 2 \approx -0.693$ for the second-order method and $z < -1.308$ for the fourth-order method.
\end{proposition}

Propositions~\ref{prop:cross}, \ref{prop:imag}, and~\ref{prop:floor} are proved in Appendix~\ref{app:cross}.
A nonlinear network can couple two real modes into a complex pair, and complex $\zeta$ reaches smaller $|R(\zeta)|$.
This removes the representability limit at the cost of a phase error.
% VERIFY: the magnitude of this phase error has not been derived.

\begin{proposition}[Second-order reformulation]\label{prop:secondorder}
Let $g = Df\, f$.
Every solution of $\dot x = f(x)$ is a solution of $\ddot x = g(x)$ with $\dot x(0) = f(x(0))$.
For any solution of $\ddot x = g(x)$, the off-manifold variable $w = \dot x - f(x)$ satisfies
\begin{equation}\label{eq:offmanifold}
\dot w = -Df(x)\, w.
\end{equation}
\end{proposition}

\begin{lemma}[Volume of the two-step recurrence]\label{lem:det}
For the recurrence $x_{j+1} = 2x_j - x_{j-1} + c_2\,\Net(x_j)$, the map $(x_j, x_{j-1}) \mapsto (x_{j+1}, x_j)$ has Jacobian determinant $1$ for every $\Net$ and every $c_2$.
\end{lemma}

Proposition~\ref{prop:secondorder} and Lemma~\ref{lem:det} are proved in Appendix~\ref{app:secondorder}.
Methods that integrate $\ddot x = g$ without resetting $\dot x$ to an estimate of $f(x)$ at each step contain the dynamics of Eq.~\eqref{eq:offmanifold}.
In the test mode, $Df = \lambda$ and $w$ grows like $e^{-\lambda t}$.
For a normal constant $Df$, the growth rate of $w$ is $-\min \operatorname{Re}\operatorname{spec}(Df)$.
% VERIFY: for non-normal, time-dependent Df the growth rate of w is not in general the negative of a Lyapunov exponent of the flow. This is an open point.

% ===============================================================
\section{Rollout error}\label{sec:rollout}
% ===============================================================

The rollout error is $\epsilon_j = \|\hat x_j - x_j\|$, where $\hat x_j$ is the predicted state and $x_j$ is the true state.

\begin{proposition}[Rollout error of one-step methods]\label{prop:rollout}
Let the method have no history-based slots.
Let $\|G_\Meth(x) - G_\Meth(x')\| \le \chi\|x - x'\|$ on a set that contains the true and predicted states.
Then
\begin{equation}\label{eq:rollout}
\epsilon_J \le \delta(\Meth)\, \frac{\chi^J - 1}{\chi - 1}, \qquad \chi \ne 1,
\end{equation}
and $\epsilon_J \le J\,\delta(\Meth)$ for $\chi = 1$.
With $\chi = e^{\gamma h}$, the amplification factor is $(e^{\gamma t_{\mathrm{end}}} - 1)/(e^{\gamma h} - 1)$.
\end{proposition}

Proposition~\ref{prop:rollout} is proved in Appendix~\ref{app:rollout}.
It is Eq.~(20) of \cite{ainslie2026} with the injected error written out.
The amplification factor is analyzed in \cite{ainslie2026}.
The injected error is bounded by Theorems~\ref{thm:affine} and~\ref{thm:nonaffine}.

For methods with history-based slots, the test-mode rollout recurrence has characteristic roots $e^z$ and $\xi$.
The amplification factor over $J$ steps is then of order $\sum_{j<J} \max(|e^z|, |\xi|)^j$.
For the Taylor method of order 2 with the history-based backward difference, $\xi = e^{-z}$.
On a decaying mode, $|\xi| > 1$.

% ===============================================================
\section{Procedure for a proposed method}\label{sec:procedure}
% ===============================================================

The following steps evaluate a proposed method.
Each step points to the result it uses.

\begin{enumerate}[label=\arabic*.]
\item Write the step map in the form of Eq.~\eqref{eq:method}.
Classify each lower slot as an exact slot, a state-based slot, or a history-based slot (Definition~\ref{def:slottype}).
\item If a lower slot is history-based, compute the parasitic roots in the test mode (Definition~\ref{def:permode}).
If $|\xi| > 1$ on modes present in the data, reject the method or reset the slot from a state-based estimate at each step.
If the method integrates $\ddot x = g$ with a tracked velocity, apply Proposition~\ref{prop:secondorder}.
\item Check representability on the resolved modes and the unresolved modes (Definition~\ref{def:permode}, Proposition~\ref{prop:floor}).
\item Compute the hardness.
Use Corollary~\ref{cor:taylor} or Theorem~\ref{thm:nonaffine} for the resolved modes.
Use Table~\ref{tab:permode} for the unresolved modes.
\item Add the slot errors.
Add $\sum_k |c_k| \varepsilon^{(k)}_\sys$ to the hardness and $\sum_k |c_k|\varepsilon^{(k)}_\ran$ to the floor (Theorem~\ref{thm:affine}).
Check Rule~\ref{rule:fd} and Eq.~\eqref{eq:noisewindow}.
\item Choose the capacity regime, which sets $\rho$ and $\alpha$ from $W$, $m$, $\nu$, and $d$ (Assumption~\ref{ass:capacity}).
\item Compare two methods $\Meth$ and $\Meth'$ in the same capacity regime by
\begin{equation}\label{eq:compare}
\frac{\epsilon_J(\Meth)}{\epsilon_J(\Meth')} \approx \frac{\Gamma(\Meth)}{\Gamma(\Meth')}\Big(\frac{\Hard(\Meth)}{\Hard(\Meth')}\Big)^{\alpha},
\end{equation}
where $\Gamma$ is the amplification factor.
Eq.~\eqref{eq:compare} holds when the floor terms are small compared with the hardness terms.
The rate $\rho$ cancels in Eq.~\eqref{eq:compare}.
The network size enters through $\alpha$ and through the size of the floor terms relative to $\rho\,\Hard^\alpha$.
\end{enumerate}

% ===============================================================
\section{Worked cases}\label{sec:cases}
% ===============================================================

\subsection{Direct step and Euler step}
By Corollary~\ref{cor:taylor}, the Euler step has hardness at most $Le^{hL}h$ and the direct step has hardness at least $e^{-hL}$.
By Theorem~\ref{thm:affine} with $\alpha = 1$, the ratio of the hardness bounds is at most $L e^{2hL} h$.
By Proposition~\ref{prop:cross}(a), a decaying mode with $z < -\ln 2$ has smaller per-mode hardness under the direct step.
The Euler step is preferred on resolved modes.
The direct step is preferred on decaying modes with $z < -\ln 2$.

\subsection{Runge--Kutta wrappers}
By Theorem~\ref{thm:nonaffine}, a Runge--Kutta wrapper of any order has leading hardness $hL$, the same as the Euler step.
The order changes the target at order $h^r$.
By Proposition~\ref{prop:floor}, the fourth-order wrapper cannot represent decaying real modes with $z < -1.308$.
By Proposition~\ref{prop:cross}(b), the implicit Euler wrapper has smaller per-mode hardness than the Euler step on every decaying mode.

\subsection{Taylor method of order 2 with exact lower slot}
By Corollary~\ref{cor:taylor}, the hardness is at most $L_2 e^{hL} h^2/2$.
By Rule~\ref{rule:order}, this bound is smaller than the Euler bound if $hL_2 < 2L$.
By Proposition~\ref{prop:cross}(g), the per-mode hardness grows like $|z|$ on decaying unresolved modes.
The exact lower slot contributes the term $z$ to the known multiplier, and the network must cancel it.
An exponential known part for the linear operator removes this growth (Table~\ref{tab:permode}).

\subsection{Taylor method of order 2 with backward-difference lower slot}
During training the slot is state-based.
By Lemma~\ref{lem:fd} with $q = 1$, $\varepsilon^{(1)}_\sys \le L_2 e^{hL} h/2$.
By Theorem~\ref{thm:affine}, the hardness term is at most $L_2 e^{hL} h^2$, twice the exact-slot bound.
During rollout the slot is history-based.
The parasitic root is $e^{-z}$ (Table~\ref{tab:permode}).
For decaying modes, $|e^{-z}| > 1$.
The method is not suitable for rollout on dissipative systems unless the slot is reset from a state-based estimate at each step.

\subsection{Runge--Kutta--Nystr\"om integration of the second-order term}\label{sec:rkn}
A Runge--Kutta--Nystr\"om method for $\ddot x = \Net(x)$ with velocity estimate $v_j$ has the form
\begin{equation}\label{eq:rkn}
x_{j+1} = x_j + h v_j + h^2 \sum_i \bar b_i\, \Net(X_i), \qquad X_i = x_j + \tilde c_i h v_j + h^2 \sum_{l<i} \bar a_{il}\, \Net(X_l),
\end{equation}
with $\sum_i \bar b_i = 1/2$.
Here $\tilde c_i$ are the stage times of the method.
By Proposition~\ref{prop:secondorder}, the true trajectory solves $\ddot x = g(x)$ with $\dot x = f(x)$.
The source of $v_j$ sets the behavior.
\begin{enumerate}[label=(\alph*)]
\item $v_j = f(x_j)$, an exact slot.
The method is a one-step method.
In the test mode, the known multiplier contains $1 + z$, and $\Hard(\Meth; z) = |e^z - 1 - z|\,(1 + O(|z|))$ on resolved modes.
The leading hardness equals that of the Taylor method of order 2 with exact lower slots.
Under Assumption~\ref{ass:target} applied to Eq.~\eqref{eq:rkn}, the target differs from $g$ by $O(h)$ or less, with the power set by the order conditions of the tableau.
% VERIFY: the power for specific RKN tableaux.
On decaying unresolved modes the per-mode hardness grows like $|z|$, as in Proposition~\ref{prop:cross}(g).
\item $v_j$ from a backward difference of predicted states, a history-based slot.
The rollout recurrence has the parasitic root $e^{-z}$.
\item $v_j$ tracked as a state variable, $v_{j+1} = v_j + h\sum_i b_i \Net(X_i)$.
The continuous system contains Eq.~\eqref{eq:offmanifold}.
Deviations from $v = f(x)$ grow on decaying modes for any integrator.
\item $v_j$ from a state-based network for $x^{(1)}$ with slot error $e_1$.
The method is a one-step method.
By Rule~\ref{rule:fd}, the order-2 hardness is kept if $\varepsilon^{(1)}_\sys = O(h)$.
The random part adds $h\,\varepsilon^{(1)}_\ran$ to the floor.
\end{enumerate}
Configurations (a) and (d) are one-step methods.
Relative to the Taylor method of order 2, they change the target, not the leading hardness.

% ===============================================================
\section{Testable predictions}\label{sec:predictions}
% ===============================================================

\begin{enumerate}[label=P\arabic*.]
\item At fixed $W$ and $m$, the injected error scales like $h^0$ for the direct step, $h^{\alpha}$ for the Euler step, and $h^{p\alpha}$ for the Taylor method of order $p$ with exact lower slots, on resolved modes (Corollary~\ref{cor:taylor}, Theorem~\ref{thm:affine}).
\item At small $h$, the Euler step and Runge--Kutta wrappers of any order have injected errors that agree to leading order (Theorem~\ref{thm:nonaffine}).
\item The injected error spectrum per wavenumber follows the per-mode hardness of Table~\ref{tab:permode} evaluated at $z_k = h\lambda_k$.
\item The Taylor method of order 2 with exact lower slots has large injected error on decaying unresolved modes. An exponential known part removes it (Proposition~\ref{prop:cross}(g)).
\item The fourth-order Runge--Kutta wrapper has an injected error floor on real decaying modes with $z < -1.308$ (Proposition~\ref{prop:floor}).
\item The Taylor method of order 2 with the history-based backward difference diverges on dissipative systems and does not diverge from the parasitic root on systems with purely imaginary test-mode rates (Table~\ref{tab:permode}).
\item Replacing an exact lower slot by a finite-difference estimate of order $q$ keeps the order-$p$ scaling of P1 if and only if $q \ge p-1$ (Rule~\ref{rule:fd}).
\item In the approximation-limited regime, the ratio of injected errors of two methods does not depend on $W$ (Eq.~\eqref{eq:compare}).
\end{enumerate}

% ===============================================================
\section{Status of the results}\label{sec:status}
% ===============================================================

Lemmas~\ref{lem:scaling}, \ref{lem:taylor}, \ref{lem:direct}, \ref{lem:fd}, \ref{lem:rk}, and~\ref{lem:det} and Propositions~\ref{prop:table}--\ref{prop:secondorder} and~\ref{prop:rollout} are proved under the stated assumptions.
Theorems~\ref{thm:affine} and~\ref{thm:nonaffine} are proved given Assumption~\ref{ass:capacity}, and Theorem~\ref{thm:nonaffine} also needs Assumptions~\ref{ass:target} and~\ref{ass:wrapper}.
Assumption~\ref{ass:capacity} is a model.
Its form with a single exponent $\alpha$ can be tested by fitting the injected error against the hardness across methods at fixed $W$ and $m$.
The test-mode calculus of Section~\ref{sec:testmode} is exact for linear normal dynamics.
For nonlinear dynamics it is a leading-order model on each mode.
A proof that joins the resolved-mode results of Section~\ref{sec:affine} with the unresolved-mode results of Section~\ref{sec:testmode} is not given here.
An inertial-manifold splitting is one candidate.

% ===============================================================
\appendix
\section{Proofs}\label{app:proofs}
% ===============================================================

\subsection{Proof of Lemma~\ref{lem:scaling}}\label{app:scaling}
Let $U = \Phi_h - Q_h^\ex$.
The step error is $G_\Meth - \Phi_h = c_p\Net - U$.
The network is trained on labels $y = (\Phi_h(x) + \eta - Q_h^\ex(x))/c_p$.
These are the function $U/c_p$ plus label noise of standard deviation $\nu/|c_p|$.
By Eq.~\eqref{eq:capacity},
\begin{equation*}
\supA{\Net - U/c_p} \le \Err\big(\Lip(U)/|c_p|,\, \nu/|c_p|\big) + \delta_{\mathrm{opt}}.
\end{equation*}
Multiply by $|c_p|$ and apply Eq.~\eqref{eq:homog} with $\vartheta = |c_p|$:
\begin{equation*}
\supA{c_p\Net - U} \le \Err\big(\Lip(U), \nu\big) + |c_p|\,\delta_{\mathrm{opt}}.
\end{equation*}
The left side is $\delta(\Meth)$. \qed

\subsection{Proof of Theorem~\ref{thm:affine}}\label{app:affine}
Write $Q_h = Q_h^\ex + \sum_k c_k e_k^\sys + \sum_k c_k e_k^\ran$.
Define $U = \Phi_h - Q_h^\ex - \sum_k c_k e_k^\sys$.
The step error is
\begin{equation*}
G_\Meth - \Phi_h = \big(c_p\Net - U\big) + \sum_k c_k e_k^\ran.
\end{equation*}
The labels are $U + \eta - \sum_k c_k e^\ran_k$, divided by $c_p$.
The noise term $\eta - \sum_k c_k e^\ran_k$ has zero mean given the state and variance $\nu_{\mathrm{eff}}^2$.
By the argument of Appendix~\ref{app:scaling}, $\supA{c_p\Net - U} \le \Err(\Lip(U), \nu_{\mathrm{eff}}) + |c_p|\delta_{\mathrm{opt}}$.
The Lipschitz seminorm is subadditive and absolutely homogeneous, so
\begin{equation*}
\Lip(U) \le \Lip(\Phi_h - Q_h^\ex) + \sum_k |c_k|\, \Lip(e_k^\sys) = \Hard(\Meth) + \sum_k |c_k|\,\varepsilon^{(k)}_\sys.
\end{equation*}
Since $\Err$ is nondecreasing in its first argument, Eq.~\eqref{eq:rate} gives the first term of Eq.~\eqref{eq:thm1}.
The triangle inequality gives $\supA{\sum_k c_k e_k^\ran} \le \sum_k |c_k|\varepsilon^{(k)}_\ran$. \qed

\subsection{Proof of Lemma~\ref{lem:taylor}}\label{app:taylor}
Fix $x \in \Aset$.
By invariance, $\Phi_t(x) \in \Aset$ for all $t$.
By Assumption~\ref{ass:regular}, $\frac{d^k}{dt^k}\Phi_t(x) = x^{(k)}(\Phi_t(x))$ for $k \le n+1$.
Taylor's theorem with integral remainder gives
\begin{equation}\label{eq:taylorrem}
\Phi_h(x) - Q^{(n)}_h(x) = \int_0^h \frac{(h-t)^n}{n!}\, x^{(n+1)}\big(\Phi_t(x)\big)\,dt.
\end{equation}
The norm of the integrand is at most $M_{n+1}(h-t)^n/n!$.
Integration gives $M_{n+1} h^{n+1}/(n+1)!$, which is Eq.~\eqref{eq:taylorsup}.

For $x, x' \in \Aset$, subtract Eq.~\eqref{eq:taylorrem} at $x'$ from Eq.~\eqref{eq:taylorrem} at $x$.
The integrand difference is bounded by $\frac{(h-t)^n}{n!} L_{n+1}\|\Phi_t(x) - \Phi_t(x')\|$.
By Gr\"onwall's inequality on $\Aset$, $\|\Phi_t(x) - \Phi_t(x')\| \le e^{tL}\|x - x'\| \le e^{hL}\|x - x'\|$ for $0 \le t \le h$.
Integration gives Eq.~\eqref{eq:taylorlip}. \qed

\subsection{Proof of Lemma~\ref{lem:direct}}\label{app:direct}
By invariance, $\Phi_{-h}$ maps $\Aset$ to $\Aset$.
By Gr\"onwall's inequality applied backward in time, $\|\Phi_{-h}(\tilde x) - \Phi_{-h}(\tilde x')\| \le e^{hL}\|\tilde x - \tilde x'\|$ for $\tilde x, \tilde x' \in \Aset$.
Set $\tilde x = \Phi_h(x)$ and $\tilde x' = \Phi_h(x')$.
Then $\|x - x'\| \le e^{hL}\|\Phi_h(x) - \Phi_h(x')\|$.
Rearranging gives the first claim.
Taking $x \ne x'$ gives $\Lip(\Phi_h) \ge e^{-hL}$.
The direct step has $Q_h^\ex = 0$, so $\Hard = \Lip(\Phi_h)$. \qed

\subsection{Proof of Lemma~\ref{lem:fd}}\label{app:fd}
For $i \ge 0$, apply Taylor's theorem to $t \mapsto \Phi_t(x)$ on the interval from $0$ to $-ih$:
\begin{equation*}
\Phi_{-ih}(x) = \sum_{k=0}^{q} \frac{(-ih)^k}{k!}\, x^{(k)}(x) + T_i(x), \qquad T_i(x) = \int_0^{-ih} \frac{(-ih - t)^q}{q!}\, x^{(q+1)}\big(\Phi_t(x)\big)\,dt.
\end{equation*}
Multiply by $\phi_i/h$ and sum over $i$.
By Eq.~\eqref{eq:fdorder}, the coefficient of $h^{k-1} x^{(k)}$ is $1$ for $k = 1$ and $0$ for $k \ne 1$.
Therefore $F_1 - f = h^{-1}\sum_i \phi_i T_i$.
For $x, x' \in \Aset$, the argument of Appendix~\ref{app:taylor} run backward in time gives
\begin{equation*}
\|T_i(x) - T_i(x')\| \le L_{q+1}\, e^{ihL}\, \frac{(ih)^{q+1}}{(q+1)!}\, \|x - x'\|.
\end{equation*}
Sum over $i$ with weights $|\phi_i|/h$ and use $e^{ihL} \le e^{qhL}$ for $i \le q$.
This gives Eq.~\eqref{eq:fd}.
For $q = 1$, $(\phi_0, \phi_1) = (1, -1)$ satisfies Eq.~\eqref{eq:fdorder}, and $C_1 = 1 \cdot 1^2/2 = 1/2$.
For $q = 2$, $(\phi_0, \phi_1, \phi_2) = (3/2, -2, 1/2)$ satisfies Eq.~\eqref{eq:fdorder}, and $C_2 = (2 \cdot 1 + \tfrac12 \cdot 8)/6 = 1$. \qed

\subsection{Proof of Lemma~\ref{lem:rk}}\label{app:rk}
Let $\kappa_i^{(1)}$ and $\kappa_i^{(2)}$ be the stages for $\mathcal N_1$ and $\mathcal N_2$, with stage points $X_i^{(1)}$ and $X_i^{(2)}$.
Set $\Delta_i = \|\kappa_i^{(1)} - \kappa_i^{(2)}\|$ and $\Delta_0 = \|\mathcal N_1 - \mathcal N_2\|_\infty$.
Then
\begin{equation*}
\Delta_i \le \|\mathcal N_1(X_i^{(1)}) - \mathcal N_1(X_i^{(2)})\| + \|\mathcal N_1(X_i^{(2)}) - \mathcal N_2(X_i^{(2)})\| \le \Lambda h \sum_{l<i} |a_{il}|\,\Delta_l + \Delta_0.
\end{equation*}
By induction, $\Delta_i \le \Delta_0(1 + h\Lambda\bar a)^{i-1}$.
The case $i = 1$ holds because the sum is empty.
If the bound holds for all $l < i$, then $\Delta_i \le \Delta_0 + h\Lambda\bar a\,\Delta_0(1 + h\Lambda\bar a)^{i-2} \le \Delta_0(1 + h\Lambda\bar a)^{i-1}$.
The step maps differ by $h\sum_i b_i(\kappa_i^{(1)} - \kappa_i^{(2)})$, with norm at most $h\beta\Delta_0(1 + h\Lambda\bar a)^{s-1}$. \qed

\subsection{Proof of Theorem~\ref{thm:nonaffine}}\label{app:nonaffine}
By Assumption~\ref{ass:target}, $G_\Meth[\Net] - \Phi_h = G_\Meth[\Net] - G_\Meth[\Netstar]$ on $\Aset$.
By Lemma~\ref{lem:rk} with $\mathcal N_1 = \Net$ and $\mathcal N_2 = \Netstar$,
\begin{equation*}
\delta(\Meth) \le h\beta(1 + h\Lambda\bar a)^{s-1}\,\|\Net - \Netstar\|_\infty.
\end{equation*}
By Assumption~\ref{ass:wrapper}, $\|\Net - \Netstar\|_\infty \le \Err(\Lip(\Netstar), \nu/h) + \delta_{\mathrm{opt}}$.
By Eq.~\eqref{eq:homog} with $\vartheta = h$,
\begin{equation*}
h\,\Err(\Lip(\Netstar), \nu/h) = \Err(h\Lip(\Netstar), \nu) = \Err(\Hard(\Meth), \nu).
\end{equation*}
This gives Eq.~\eqref{eq:thm2}.
By Assumption~\ref{ass:target}, $\Lip(\Netstar) = L + O(h^r)$. \qed

\subsection{Proof of Proposition~\ref{prop:table}}\label{app:table}
The exact multiplier of the test mode over one step is $e^z$.
\begin{enumerate}[label=(\roman*)]
\item Direct step. $R = \mu$, $\mu^\star = e^z$, $\partial_\mu R = 1$, so $\Hard = |e^z|$.
\item Euler step. $R = 1 + h\mu$, $h\mu^\star = e^z - 1$, $\partial_\mu R = h$, so $\Hard = |e^z - 1|$.
\item Implicit Euler wrapper, $x_{j+1} = x_j + h\mu x_{j+1}$. $R = 1/(1 - h\mu)$. The condition $R = e^z$ gives $h\mu^\star = 1 - e^{-z}$, which exists for all $z$. $\partial_\mu R = h/(1 - h\mu)^2 = h e^{2z}$ at $\mu^\star$. So $\Hard = |1 - e^{-z}|\,|e^{2z}| = |e^z|\,|e^z - 1|$.
\item Explicit Runge--Kutta wrapper. $R_\Meth = R(h\mu)$ with $R$ the stability polynomial. Representability requires $e^z \in R(\C)$. With $\zeta = h\mu^\star$, $\partial_\mu R_\Meth = hR'(\zeta)$, so $\Hard = |\zeta R'(\zeta)|$.
\item Taylor method of order $p$ with exact lower slots. $x^{(k)}(x) = \lambda^k x$, so $\psi(z) = \sum_{k<p} z^k/k! = \expt{p-1}(z)$. $\Hard = |e^z - \expt{p-1}(z)|$.
\item Taylor method of order 2 with backward difference. On the true trajectory, $x_{j-1} = e^{-z}x_j$, so $hF_1 = (1 - e^{-z})x_j$ and $\psi(z) = 2 - e^{-z}$. $\Hard = |e^z - 2 + e^{-z}|$. In rollout, $hF_1 = x_j - x_{j-1}$ and the network term equals $(e^z - 2 + e^{-z})x_j$ at $\mu^\star$. The recurrence is $x_{j+1} = (e^z + e^{-z})x_j - x_{j-1}$. Its characteristic polynomial is $\xi^2 - (e^z + e^{-z})\xi + 1 = (\xi - e^z)(\xi - e^{-z})$. The parasitic root is $e^{-z}$.
\item Exponential known part. $\psi(z) = e^{z_{\mathrm{lin}}}$, so $\Hard = |e^z - e^{z_{\mathrm{lin}}}| = |e^{z_{\mathrm{lin}}}|\,|e^{z - z_{\mathrm{lin}}} - 1|$. \qed
\end{enumerate}

\subsection{Proofs of Propositions~\ref{prop:cross}, \ref{prop:imag}, and~\ref{prop:floor}}\label{app:cross}
All $z$ are real in Proposition~\ref{prop:cross}.
\begin{enumerate}[label=(\alph*)]
\item For $z \ge 0$, $e^z - 1 < e^z$. For $z < 0$, $1 - e^z < e^z$ if and only if $e^z > 1/2$, that is $z > -\ln 2$.
\item The ratio of the implicit Euler and Euler per-mode hardness is $e^z$, which is below $1$ if and only if $z < 0$.
\item The ratio of the implicit Euler and direct-step per-mode hardness is $|e^z - 1|$, which is below $1$ if and only if $z < \ln 2$.
\item $e^z - 1 - z \ge 0$ for all real $z$. The condition $e^z - 1 - z < e^z$ is $z > -1$.
\item For $z \ge 0$, $e^z - 1 - z < e^z - 1$. For $z < 0$, the condition is $e^z - 1 - z < 1 - e^z$, that is $2e^z - 2 - z < 0$. The function $2e^z - 2 - z$ is convex, is zero at $z = 0$ and at $z_\ast$, and is negative between them. Numerical root finding gives $z_\ast = -1.5936$.
\item Set $u = e^z$. The per-mode hardness is $u + 1/u - 2 \ge 0$. Against the direct step, $u + 1/u - 2 < u$ if and only if $u > 1/2$. Against the Euler step with $u < 1$, $u + 1/u - 2 < 1 - u$ is $2u^2 - 3u + 1 < 0$, that is $1/2 < u < 1$. Against the Euler step with $u \ge 1$, $u + 1/u - 2 < u - 1$ is $1/u < 1$, which holds for $u > 1$. Both conditions reduce to $z > -\ln 2$.
\item For $z < 0$, $|e^z - 1 - z| = e^z - 1 - z = |z| - 1 + e^z$.
\end{enumerate}
For Proposition~\ref{prop:imag}, $|e^{\mathrm{i}\omega}| = 1$ and $|e^{\mathrm{i}\omega} - 1| = 2|\sin(\omega/2)|$.
The expansion $e^{\mathrm{i}\omega} - 1 - \mathrm{i}\omega = -\omega^2/2 + O(\omega^3)$ gives the order-2 value.
Also $e^{\mathrm{i}\omega} + e^{-\mathrm{i}\omega} - 2 = 2\cos\omega - 2 = -4\sin^2(\omega/2)$.

For Proposition~\ref{prop:floor}, the second-order polynomial $1 + \zeta + \zeta^2/2$ has derivative $1 + \zeta$, zero at $\zeta = -1$, with value $1/2$.
The fourth-order polynomial $\expt{4}(\zeta)$ has derivative $\expt{3}(\zeta)$.
This cubic has one real root, $\zeta \approx -1.5961$, found numerically.
The value there is $0.2704$.
Both polynomials tend to $+\infty$ as $\zeta \to \pm\infty$, so these values are the minima on the real line.
The limits on $z$ are $\ln(1/2) = -0.693$ and $\ln(0.2704) = -1.308$. \qed

\subsection{Proofs of Proposition~\ref{prop:secondorder} and Lemma~\ref{lem:det}}\label{app:secondorder}
If $\dot x = f(x)$, then $\ddot x = Df(x)\dot x = Df(x) f(x) = g(x)$ and $\dot x(0) = f(x(0))$.
For any solution of $\ddot x = g(x)$, set $w = \dot x - f(x)$.
Then $\dot w = \ddot x - Df(x)\dot x = Df(x)f(x) - Df(x)\dot x = -Df(x)\,w$.

For Lemma~\ref{lem:det}, the Jacobian of $(x_j, x_{j-1}) \mapsto (x_{j+1}, x_j)$ is the block matrix
\begin{equation*}
\begin{pmatrix} 2I + c_2\, D\Net(x_j) & -I \\ I & 0 \end{pmatrix}.
\end{equation*}
For a block matrix $\begin{pmatrix} Y_{11} & Y_{12} \\ Y_{21} & 0\end{pmatrix}$ with $n_{\mathrm{x}} \times n_{\mathrm{x}}$ blocks, exchanging the two block columns gives the determinant $(-1)^{n_{\mathrm{x}}}\det(Y_{12})\det(Y_{21})$.
Here $\det(Y_{12}) = (-1)^{n_{\mathrm{x}}}$ and $\det(Y_{21}) = 1$, so the determinant is $1$. \qed

\subsection{Proof of Proposition~\ref{prop:rollout}}\label{app:rollout}
By the triangle inequality,
\begin{equation*}
\epsilon_{j+1} = \|G_\Meth(\hat x_j) - \Phi_h(x_j)\| \le \|G_\Meth(\hat x_j) - G_\Meth(x_j)\| + \|G_\Meth(x_j) - \Phi_h(x_j)\|.
\end{equation*}
The first term is at most $\chi\,\epsilon_j$.
The second term is at most $\delta(\Meth)$ because $x_j \in \Aset$.
With $\epsilon_0 = 0$, induction gives $\epsilon_J \le \delta(\Meth)\sum_{j=0}^{J-1} \chi^j$.
The geometric sum gives Eq.~\eqref{eq:rollout}. \qed

% ===============================================================
\begin{thebibliography}{9}
\bibitem{ainslie2026}
C.~Ainslie, P.~Hassanzadeh, M.~W.~Mahoney, and A.~Chattopadhyay.
Eigenanalysis framework for autoregressive neural emulators of multi-scale chaotic dynamics.
arXiv:2608.16084, 2026.

\bibitem{mohri2008}
M.~Mohri and A.~Rostamizadeh.
Rademacher complexity bounds for non-i.i.d.\ processes.
In \emph{Advances in Neural Information Processing Systems 21}, 2008.

\bibitem{stone1982}
C.~J.~Stone.
Optimal global rates of convergence for nonparametric regression.
\emph{Annals of Statistics}, 10(4):1040--1053, 1982.
% VERIFY: volume and pages.

\bibitem{yarotsky2017}
D.~Yarotsky.
Error bounds for approximations with deep ReLU networks.
\emph{Neural Networks}, 94:103--114, 2017.
% VERIFY: volume and pages.

\bibitem{zhu2020imde}
A.~Zhu, P.~Jin, and Y.~Tang.
Inverse modified differential equations for discovery of dynamics.
arXiv:2009.01058, 2020.
\end{thebibliography}

\end{document}
